\section{Experimental design}
\label{sec:protocols}

\begin{center}
\small
\begin{tabularx}{\textwidth}{@{}l
  >{\raggedright\arraybackslash}p{0.22\textwidth}
  >{\raggedright\arraybackslash}p{0.25\textwidth}
  >{\raggedright\arraybackslash}X@{}}
\toprule
Stage & Model \& parameter count & Core questions & Pre-specified outcome \\
\midrule
Phase~1 & Coupled double well and shared-head MLP (\PoneDimensions{} parameters) &
Native replay, energy, known-boundary recovery, and continual retention &
Numerical gates and finite-horizon recovery passed. Retention failed. \\
\addlinespace
Phase~2 & Shared-head MLP (\PoneDimensions{} parameters) &
Replay precision, delayed feedback, and Fisher-based retention &
Integer replay passed on the tested grid. Terminal recovery passed, but ramped recovery and retention failed. \\
\addlinespace
Phase~3 & Causal Transformer (\PthreeParams{} parameters) &
Lattice inversion, search under terminal feedback, and checkpoint costs &
The combined precision claim failed despite exact tree replay. Search and the systems criteria passed. \\
\bottomrule
\end{tabularx}
\end{center}

For each confirmatory stage, we fixed the data generator, comparisons,
outcomes, and decision rules before examining held-out results. Separate
adapter implementations used the same evaluators. Where we tuned
settings, validation and test seeds were disjoint, and we locked the
adapter with its validation receipt before the held-out run. A seed
defines one model and dataset, so rerunning it on another backend checks
reproducibility of the same replicate. Appendix~\ref{app:reproducibility}
lists the underlying artifacts and commands.

\paragraph{Analytic and neural precision.}
We measured energy error across step sizes in the Phase~1 double well
with the potential held fixed. A separate tilted tape tested momentum
reflection and reverse event order under changing forces, as did the
neural tape on a shared-head multilayer perceptron (MLP). In Phase~2
we compared native binary32 endpoints with a fixed-point control on
a pre-specified replay grid. We ran a separate dense scan to locate the
first tolerance miss between grid points.

\paragraph{Feedback and retention.}
The learning tasks classify procedural clusters whose centers and
context change between tasks. The poison and retention experiments
start from the same first-task pretraining state and run on separate
branches. In recovery, the protocol specifies when a policy gets
feedback, how far it can rewind, and when the scoring data are used.
In retention, all methods carry their state across tasks and train
without replaying batches from earlier tasks. If \(a_{i,j}\) is accuracy
on task \(j\) after training through task \(i\), we measure retention
as \(\BWT=a_{J,1}-a_{1,1}\) at the last task \(J\), and live accuracy
as the mean of \(a_{j,j}\) on the new tasks. We pair BWT contrasts
by test seed. The pre-specified sign tests compare each contrast with a
positive superiority margin and exclude ties at that margin, computed
on the integer correct-count scale.

Phase~1 matches arriving minibatches at fixed rates. In Phases~2
and~3 we count fresh data gradients and charge separately for each replay
gradient. The Phase~2 retention settings come from validation, and all
methods share a continual gradient budget. In recovery, policies spend
any remaining credits on the same ordered stream of clean repair
batches. Per-batch scoring and data transfers add work outside the
gradient count.

\paragraph{Speculative search.}
In Phase~3 we train all \PthreeParams{} parameters of the Transformer
in the full-model runs. Each procedural problem has \(D\) public
addresses with default continuation bits. The private rule is drawn
uniformly from the default and its \(D\) single-address exceptions.
Training examples are clean conditional on the proposed rule. Search
follows a fixed depth-first trie, while the verifier keeps the private
rule and checks held-out model predictions. Since the policy responds
to a verifier outcome only by stopping at success, a single exhaustive
traversal provides the leaf predictions and stopping prefixes for every
possible private rule.

Reversible methods save ordered event IDs and rewind the relevant
suffixes. Checkpoint methods restore weights and optimizer state at
unresolved forks from GPU buffers or pinned host memory. The single-path
control submits the first native leaf to the verifier once. In separate
correctness traversals we save parent snapshots to check every unwound
edge. For timings we omit those snapshots, run each method and depth in
a fresh CUDA process, and synchronize the device for block timers. We
measure peak allocation in those processes.
